\section{结论}
\label{sec:conclusion}

本文提出了 PivotRL，这是一种轮次级强化学习算法，旨在以可承受的方式为大语言模型执行长时程智能体任务的后训练。我们在四个不同的智能体领域上的实验表明，该方法同时具备有效性与可扩展性。与使用相同数据的标准 SFT 相比，PivotRL 的域内平均准确率高出 $+4.17$，同时在 OOD 保留上多出 $+10.04$ 的优势，显著缓解了灾难性遗忘。进一步地，与端到端 RL 相比，PivotRL 能以更低得多的计算成本取得相近准确率。PivotRL 的性能收益主要来自两点：一是 on-policy rollout 带来了更充分的动作空间覆盖，二是 verifier 产生的奖励能为功能正确的动作分配信用。未来，我们计划把该框架扩展到非程序化 verifier，例如 LLM-as-a-judge 框架 \citep{zheng2023judgingllmasajudgemtbenchchatbot} 和过程奖励模型 \citep{lightman2023letsverifystepstep}，同时探索动态采样等在线奖励剖析方法 \citep{yu2025dapoopensourcellmreinforcement}。
